\section{平坦模与``Tor''函子}

为了方便熟悉同调代数 (*) 的读者，我们简要指出平坦模的理论与 Tor 函子的理论之间的联系。

命题 1. 设 E 是右 A-模。下列四个性质等价：

\begin{enumerate}
\def\labelenumi{(\alph{enumi})}
\item
  E 是平坦的。
\item
  对每个左 A-模 F 与每个整数 \(n \geqslant 1\) ， \(\operatorname{Tor}_n^{\mathbf{A}}(\mathbf{E}, \mathbf{F}) = 0\) 。
\item
  对每个左 A-模 F， \(\operatorname{Tor}_{1}^{\mathbf{A}}(\mathbf{E},\mathbf{F}) = 0\) 。
\item
  对 A 的每个有限生成左理想 a，
\end{enumerate}

\[
\operatorname{Tor} _ {1} ^ {\mathbf {A}} (\mathrm{E}, \mathrm{A} _ {s} / \mathfrak {a}) = 0.
\]

我们证明 (a) 蕴含 (b)。设

\[
\dots \rightarrow \mathrm{L} _ {n} \rightarrow \mathrm{L} _ {n - 1} \rightarrow \dots \rightarrow \mathrm{L} _ {0} \rightarrow \mathrm{F} \rightarrow 0
\]

是 F 的自由分解。由于 E 平坦，序列

\[
\dots \rightarrow \mathrm{E} \otimes \mathrm{L} _ {n} \rightarrow \mathrm{E} \otimes \mathrm{L} _ {n - 1} \rightarrow \dots \rightarrow \mathrm{E} \otimes \mathrm{L} _ {0} \rightarrow \mathrm{E} \otimes \mathrm{F} \rightarrow 0\tag{1}
\]

是正合的。由于诸 \(\operatorname{Tor}_n^{\mathbf{A}}(\mathbf{E},\mathbf{F})\) 同构于复形 (1) 的同调群，它们在 \(n\geqslant 1\) 时为零。(b) 蕴含 (c) 且 (c) 蕴含 (d) 是平凡的。最后我们证明 (d) 蕴含 (a)。正合序列

\[
0 \rightarrow a \rightarrow A _ {s} \rightarrow A _ {s} / a \rightarrow 0
\]

给出正合序列

\[
\operatorname{Tor} _ {1} ^ {\mathbf {A}} (\mathrm{E}, \mathrm{A} _ {s} / \mathfrak {a}) \rightarrow \mathrm{E} \otimes_ {\mathbf {A}} \mathfrak {a} \rightarrow \mathrm{E} \otimes_ {\mathbf {A}} \mathrm{A}.
\]

由于 (d) 成立，典范同态

\[
\mathbf {E} \otimes_ {\mathbf {A}} \mathfrak {a} \rightarrow \mathbf {E} \otimes_ {\mathbf {A}} \mathbf {A} = \mathbf {E}
\]

是单射，这就是说 E 是平坦的（§2 第 3 目，命题 1）。

命题 1 给出了平坦模的一个在应用中常常有用的判别准则。作为例子，我们只限于给出 §2 第 5 目命题 5 的一个新证明。若 E' 与 E'' 都平坦，正合序列

\[
\operatorname{Tor} _ {1} ^ {\mathbf {A}} \left(\mathrm{E} ^ {\prime}, \mathrm{F}\right)\rightarrow \operatorname{Tor} _ {1} ^ {\mathbf {A}} (\mathrm{E}, \mathrm{F}) \rightarrow \operatorname{Tor} _ {1} ^ {\mathbf {A}} \left(\mathrm{E} ^ {\prime \prime}, \mathrm{F}\right)
\]

表明 \(\operatorname{Tor}_1^{\mathbf{A}}(\mathbf{E},\mathbf{F}) = 0\) 对每个左 A-模 \(\mathbf{F}\) 成立，因此 \(\mathbf{E}\) 是平坦的。若 \(\mathbf{E}\) 与 \(\mathbf{E}''\) 都平坦，正合序列

\[
\operatorname{Tor} _ {2} ^ {\mathbf {A}} \left(\mathrm{E} ^ {\prime \prime}, \mathrm{F}\right)\rightarrow \operatorname{Tor} _ {1} ^ {\mathbf {A}} \left(\mathrm{E} ^ {\prime}, \mathrm{F}\right)\rightarrow \operatorname{Tor} _ {1} ^ {\mathbf {A}} (\mathrm{E}, \mathrm{F})
\]

表明 \(\operatorname{Tor}_1^{\mathbf{A}}(\mathbf{E}',\mathbf{F}) = 0\) ，因此 \(\mathbf{E}'\) 是平坦的。

命题 2. 设 R、S 是两个环， \(\rho: \mathbf{R} \to \mathbf{S}\) 是一个同态，F 是左 R-模。下列两个性质等价：

\begin{enumerate}
\def\labelenumi{(\alph{enumi})}
\item
  \(\operatorname{Tor}_1^{\mathbf{R}}(\rho_*(\mathbf{E}),\mathbf{F}) = 0\) ，对每个右 S-模 E。
\item
  左 S-模 \(\rho^{*}(\mathbf{F}) = \mathbf{F}_{(\mathbf{S})} = \mathbf{S} \otimes_{\mathbb{R}} \mathbf{F}\) 是平坦的，且 \(\operatorname{Tor}_{1}^{\mathbb{R}}(\rho_{*}(\mathbf{S}_{d}), \mathbf{F}) = 0\) 。
\end{enumerate}

假设 (a) 成立。取 \(E = S_{d}\) ，我们看到 \(\operatorname{Tor}_{1}^{\mathbb{R}}(\rho_{*}(S_{d}), \mathbf{F}) = 0\) 。我们还证明 \(F_{(S)}\) 是平坦 S-模。为此我们注意：若 E 是右 S-模，加法群 \(E \otimes_{S} F_{(S)}\) 等同于 \(\rho_{*}(E) \otimes_{R} F\) 。于是，若有右 S-模的正合序列

\[
0 \rightarrow \mathrm{E} ^ {\prime} \rightarrow \mathrm{E} \rightarrow \mathrm{E} ^ {\prime \prime} \rightarrow 0
\]

则利用 (a) 得到正合序列

\[
0 \rightarrow \rho_ {*} (\mathrm{E} ^ {\prime}) \otimes_ {\mathrm{R}} \mathrm{F} \rightarrow \rho_ {*} (\mathrm{E}) \otimes_ {\mathrm{R}} \mathrm{F} \rightarrow \rho_ {*} (\mathrm{E} ^ {\prime \prime}) \otimes_ {\mathrm{R}} \mathrm{F} \rightarrow 0
\]

或写作

\[
0 \rightarrow \mathrm{E} ^ {\prime} \otimes_ {\mathbf {S}} \mathrm{F} _ {(\mathbf {S})} \rightarrow \mathrm{E} \otimes_ {\mathbf {S}} \mathrm{F} _ {(\mathbf {S})} \rightarrow \mathrm{E} ^ {\prime \prime} \otimes_ {\mathbf {S}} \mathrm{F} _ {(\mathbf {S})} \rightarrow 0
\]

这就证明 \(\mathbf{F}_{(\mathbf{S})}\) 是平坦的。

反之，若 (b) 成立，首先，对每个自由右 S-模 \(L = S_{d}^{(1)}\) ，有 \(\operatorname{Tor}_{1}^{\mathrm{R}}(\rho_{*}(\mathrm{L}), \mathrm{F}) = (\operatorname{Tor}_{1}^{\mathrm{R}}(\rho_{*}(\mathrm{S}_{d}), \mathrm{F}))^{(1)} = 0\) 。每个右 S-模 E 都可写成 E = L/H 的形式，其中 L 是适当的自由 S-模；于是有正合序列

\[
0 = \operatorname{Tor} _ {1} ^ {\mathrm{R}} (\rho_ {*} (\mathrm{L}), \mathrm{F}) \rightarrow \operatorname{Tor} _ {1} ^ {\mathrm{R}} (\rho_ {*} (\mathrm{E}), \mathrm{F}) \rightarrow \rho_ {*} (\mathrm{H}) \otimes_ {\mathrm{R}} \mathrm{F} \rightarrow \rho_ {*} (\mathrm{L}) \otimes_ {\mathrm{R}} \mathrm{F}. \tag {2}
\]

但由于 \(\mathbf{F}_{(\mathbb{S})}\) 平坦，同态 \(\mathbf{H} \otimes_{\mathbb{S}} \mathbf{F}_{(\mathbb{S})} \to \mathbf{L} \otimes_{\mathbb{S}} \mathbf{F}_{(\mathbb{S})}\) 是单射，且它等同于同态

\[
\rho_ {*} (\mathrm{H}) \otimes_ {\mathbf {R}} \mathrm{F} \rightarrow \rho_ {*} (\mathrm{L}) \otimes_ {\mathbf {R}} \mathrm{F}.
\]

于是由 (2) 推出 \(\mathrm{Tor}_1^{\mathbb{R}}(\rho_*(\mathbf{E}),\mathbf{F}) = 0\) 。

注. 命题 2 也可以由下列正合序列的存在性得出

\[
\mathrm{E} \otimes_ {\mathbf {S}} \operatorname{Tor} _ {1} ^ {\mathrm{R}} (\rho_ {*} (\mathrm{S} _ {d}), \mathrm{F}) \rightarrow \operatorname{Tor} _ {1} ^ {\mathrm{R}} (\rho_ {*} (\mathrm{E}), \mathrm{F}) \rightarrow \operatorname{Tor} _ {1} ^ {\mathrm{S}} (\mathrm{E}, \mathrm{S} _ {d} \otimes_ {\mathrm{R}} \mathrm{F}) \rightarrow 0
\]

它由 Tor 函子的``结合性''谱序列导出。

